% ======================================================================
% Section 1 -- Introduction (~1.25 pp incl. Fig 1)
% Binding: "what explainability buys practitioners"; "neighborhoods",
% never "relations"; no em dashes; short sentences; molecule gloss
% appears once, TopoTune parenthetical style.
% ======================================================================
\section{Introduction}
\label{sec:introduction}

Graph neural networks (GNNs) learn from pairwise connections. Many
systems carry structure beyond pairs: a benzene ring is a six-atom
object whose meaning is lost when read bond by bond. Topological deep
learning (TDL) models this higher-order structure directly
\citep{hajij2023tdl,papamarkou2024position}. Its domains are
\emph{combinatorial complexes}, which store multi-way interactions as
first-class objects called \emph{cells}, organized by \emph{rank}: in a
molecule, atoms (nodes, i.e., cells of rank zero), bonds (edges, i.e.,
cells of rank one), and rings (cells of rank two). Topological neural
networks (TNNs) exchange messages between cells along
\emph{neighborhoods} \citep{bodnar2021cwn,papillon2025topotune,
hopse2025}, and now match or exceed GNNs on molecular and geometric
benchmarks \citep{telyatnikov2024topobench}.

These models cannot yet be explained in their own vocabulary. Post-hoc
explainers for GNNs \citep{ying2019gnnexplainer,luo2020parameterized,
yuan2021subgraphx,muschalik2024graphshapiq} operate on nodes and
edges. A ring is not expressible in their output space, so the
structures a TNN reasons with cannot appear in its explanation. A
practitioner deploying a TNN today cannot answer two basic questions:
which cells drove this prediction, and which neighborhoods drive this
model's performance.

We introduce \method, which answers both questions with one cooperative
game (\cref{fig:overview}). The players are the cells of the input
complex. The value of a coalition is the model's output when only those
cells participate: features, neighborhood matrices, and readout pooling
are all restricted to the coalition. From this game \method returns two
outputs. The \emph{attribution} assigns each cell a Shapley value
\citep{shapley1953value}, inheriting the classical axioms. The
\emph{explanation} is the smallest sub-complex sufficient to preserve
the prediction, found by greedy pruning.

The game makes no assumption about rank. On graphs, which are rank-1
complexes, it reduces to induced-subgraph removal, and the explanation
output solves the same search problem as SubgraphX
\citep{yuan2021subgraphx}. This reduction lets us validate \method on
the community's own benchmarks before using its generality: on GraphXAI
\citep{agarwal2023graphxai}, \method matches or outperforms GNN
explainers on identical subjects, splits, and metrics
(\cref{tab:graph-benchmarks}), at up to 500 times lower cost than
search-based methods. The same implementation then runs unchanged on
higher-order models, where the baselines are undefined. It recovers
motifs on lifted molecular benchmarks (\cref{tab:lifted}), is the most
faithful method on TopoBench datasets that carry no explanation ground
truth (\cref{tab:nogt}), and separates how models decide on native
triangulations (\cref{sec:predictions}).

The same game template also explains performance
(\cref{sec:performance}). With neighborhoods as players and validation
accuracy as the value, one short training run prices every
message-passing route. The resulting selection recipe costs roughly one
tenth of an architecture sweep and concedes at most two accuracy
points to it.

Explaining these models also taught us something about their inputs.
Standard feature liftings compute higher-rank features by aggregating
lower-rank ones, which caches structural facts into feature vectors.
Models trained on such caches read them instead of computing structure,
and their explanations degrade measurably. A structure-only lifting
removes the cache and improves accuracy and explanation quality
together (\cref{sec:predictions}). This is a design rule practitioners
can apply directly.

\textbf{Contributions.}
\begin{itemize}
  \item \textbf{Method.} The participation game on combinatorial
    complexes, with two outputs: per-cell Shapley attributions and the
    smallest sufficient sub-complex (\cref{sec:framework}).
  \item \textbf{Reduction.} At rank 1 the game recovers the search
    problem of SubgraphX; the remaining design choices are ablated
    empirically (\cref{sec:predictions}).
  \item \textbf{Evidence.} Fair, same-subject comparisons on GraphXAI;
    motif recovery, ground-truth-free faithfulness, and model
    diagnostics on higher-order domains; all with multi-seed error
    bars (\cref{sec:predictions}).
  \item \textbf{Architecture selection.} The performance
    instantiation and its selection recipe, priced honestly
    (\cref{sec:performance}).
  \item \textbf{The lifting finding.} Feature caches reduce
    explainability; a structure-only lifting improves accuracy and
    explanations together (\cref{sec:predictions}).
  \item \textbf{Open source.} One reproducible pipeline covering every
    number and figure in the paper.
\end{itemize}

\begin{figure}[t]
  \centering
  \includegraphics[width=\linewidth]{../figures/out/overview.pdf}
  \caption{\textbf{\method: one game, two instantiations.}
  \textbf{A.}~Players are the cells of the input complex
  $\mathcal{C}$. A coalition $S$ keeps a subset of cells;
  $v(S)$ is the model output on the induced sub-complex
  $\mathcal{C}[S]$. Averaging marginal contributions over coalitions
  yields one Shapley value $\phi_\sigma$ per cell (output 1). Greedy
  pruning yields the smallest sufficient sub-complex $S^{*}_{b}$
  (output 2). \textbf{B.}~Players are the neighborhoods
  $\mathcal{N}$ of the architecture, drawn as source rank
  $\rightarrow$ destination rank, and $v(S)$ is the validation
  accuracy of a briefly trained model using only the neighborhoods in
  $S$; the same averaging attributes performance per neighborhood
  ($\phi_{\mathcal{N}}$).}
  \label{fig:overview}
\end{figure}
